\section{Exercise 5: Trinomial Model}

\subsection{Problem Statement}
We consider a discrete-time market model where the return of the risky asset, defined as $R_t := \frac{S_t - S_{t-1}}{S_{t-1}}$, can take one of three possible values at each time step: $a$, $b$, or $c$, such that $-1 < a < b < c$. 
The returns are independent and identically distributed (i.i.d.) with strictly positive probabilities $p_a, p_b, p_c > 0$.
This market consists of one riskless asset (with constant interest rate $r$) \& one risky asset.

\subsection{Part 1: Arbitrage-Free Condition (a < r < c)}

\textbf{Question:} Show that if the market is arbitrage-free, then $a < r < c$.

\textbf{Step-by-Step Proof:}
We will prove this by contradiction. We will show that if $r \ge c$, we can construct an arbitrage opportunity. A symmetric argument applies if $r \le a$.

\vspace{0.5em}
\noindent \textit{Case 1: Assume $r \ge c$}
\begin{enumerate}
    \item \textbf{Initial Strategy ($t=0$):} We short one unit of the risky asset. This gives us $+S_0$ in cash. We take this $S_0$ and invest it entirely in the riskless bank account. 
    \item \textbf{Initial Cost:} The initial value of this portfolio is exactly $0$.
    \item \textbf{Value at Maturity ($t=T$):} 
    \begin{itemize}
        \item The money in the bank has grown to $S_0(1+r)^T$.
        \item We must buy back the shorted asset, which costs $-S_T$.
    \end{itemize}
    \item \textbf{Analyzing the Payoff:} The final value is $V_T = S_0(1+r)^T - S_T$.
    \item We know that the maximum possible return in any single period is $c$. Therefore, the absolute maximum price the stock can reach at time $T$ is if it returns $c$ every single period: $S_T \le S_0(1+c)^T$.
    \item Substituting this maximum stock price into our payoff equation gives the worst-case scenario for our portfolio:
    $$V_T \ge S_0(1+r)^T - S_0(1+c)^T$$
    \item Since we assumed $r \ge c$, it follows that $(1+r)^T \ge (1+c)^T$. Therefore:
    $$S_0(1+r)^T - S_0(1+c)^T \ge 0 \implies V_T \ge 0$$
    Our portfolio will never lose money.
    \item \textbf{Strict Positivity:} Because $b < c$ and the probability of outcome $b$ is strictly positive ($p_b > 0$), there is a non-zero probability that the stock price at $T$ will be strictly less than $S_0(1+c)^T$. In these states, $V_T > 0$.
    \item \textbf{Conclusion for Case 1:} A portfolio costing zero today that never loses money and has a positive probability of making money is an arbitrage opportunity. To avoid this, we must have $r < c$.
\end{enumerate}

\vspace{0.5em}
\noindent \textit{Case 2: Assume $r \le a$}
\begin{enumerate}
    \item \textbf{Strategy:} We do the exact opposite. We borrow money from the bank (short the riskless asset) to buy the risky asset. 
    \item \textbf{Analyzing the Payoff:} At $t=T$, the worst possible return for the stock is $a$. The lowest possible stock price is $S_0(1+a)^T$.
    \item The final value is $V_T = S_T - S_0(1+r)^T$.
    \item In the worst-case scenario, $V_T \ge S_0(1+a)^T - S_0(1+r)^T$. Since $r \le a$, this payoff is always $\ge 0$. 
    \item Again, because $p_b > 0$ and $b > a$, there is a positive probability of making a strict profit. This is an arbitrage.
    \item \textbf{Conclusion for Case 2:} To avoid this arbitrage, we must have $r > a$.
\end{enumerate}

Combining both cases, an arbitrage-free market requires $a < r < c$.

\subsection{Part 2: Market Completeness}

\textbf{Question:} Is the market complete?

\textbf{Step-by-Step Proof:}
\begin{enumerate}
    \item A market is complete if and only if every possible payoff can be perfectly replicated by a self-financing portfolio of the available traded assets.
    \item Let's look at a single time step, $T=1$. At time $t=1$, there are exactly 3 possible states of the world for the stock price:
    \begin{itemize}
        \item State 1: $S_1 = S_0(1+a)$
        \item State 2: $S_1 = S_0(1+b)$
        \item State 3: $S_1 = S_0(1+c)$
    \end{itemize}
    \item Because there are 3 states, any arbitrary derivative payoff $H$ at $t=1$ can be represented as a 3-dimensional vector in $\mathbb{R}^3$: $H = (H_a, H_b, H_c)^T$.
    \item To replicate this payoff, we hold $\Delta$ units of the risky asset and $B$ units of the riskless asset. Our portfolio's payoff vector across the 3 states will be:
    $$ \Delta \begin{pmatrix} S_0(1+a) \\ S_0(1+b) \\ S_0(1+c) \end{pmatrix} + B \begin{pmatrix} 1+r \\ 1+r \\ 1+r \end{pmatrix} $$
    \item This equation describes a linear combination of two fixed vectors. The set of all possible linear combinations of two vectors forms a 2-dimensional plane (a subspace) within the 3-dimensional space $\mathbb{R}^3$.
    \item Since a 2-dimensional plane cannot cover the entirety of a 3-dimensional space, there exist payoff vectors $H \in \mathbb{R}^3$ that simply do not lie on this plane. 
    \item Because there are payoffs that cannot be generated by any combination of $\Delta$ and $B$, these payoffs are not replicable. 
    \item Therefore, the trinomial market is \textbf{incomplete}.
\end{enumerate}